\section{Quantitative stability of the cone}

Fix $p>0$, write $r=r_n$, $a=\log n$, and work on one event of
probability at least $1-Cn^{-p}$. Proposition~\ref{prop:core-coarse},
Lemma~\ref{lem:core-green}, Freedman's inequality and a union bound give,
simultaneously for all lattice sites $|y|\le3n$,
\begin{equation}\label{eq:rate-event}
 D_n\subset B(0,Cr),\quad \max\ell_n\le Cr,\quad
 |M_n^y|\le C(\sqrt{aW_n(y)}+a),\quad
 \ell_n(y)=U_{D_n}(y)-M_n^y+O(a).
\end{equation}
Here and below the constants may be increased by increasing the exponent
in the constituent probability estimates. Cell interpolation gives
\begin{equation}\label{eq:rate-crude}
 \|\ell_n-U_{D_n}\|_\infty\le\delta,
 \qquad \delta=
 \begin{cases}C\sqrt r\,a,&d=2,\\ C\sqrt{ra},&d\ge3.
 \end{cases}
\end{equation}
For points outside $B(0,2n)$ this follows directly from the decay of the
potential, since $\ell_n=0$ there. Define
\[
 b=\inf_{y\notin D_n}\rho(y),\quad E=D_n\setminus bK,
 \quad H=U_{D_n}(y_0),\quad I=\int_E(\rho(v)-b)\,dv,
\]
where $y_0\in\partial(bK)$ is in the closure of an empty cell $C_z$.
Thus $\ell_n(z)=0$, $|z-y_0|\le\sqrt d/2$, and convexity gives
$H=U_E(y_0)\ge0$. Set
\begin{equation}\label{eq:rate-hq}
 h=\begin{cases}\sqrt{ra},&d=2,\\a,&d\ge3,\end{cases}
 \qquad q=h/r.
\end{equation}

\begin{lemma}[Noise at the first hole]\label{lem:rate-contact}
On the preceding event, $H\le Ch$.
\end{lemma}
\begin{proof}
The equation at the empty site and the cell modulus give
\begin{equation}\label{eq:rate-contact-noise}
 H\le C(\sqrt{aW_n(z)}+a).
\end{equation}
For every nonnegative integrable radial function $\chi$, Newton's shell
identity gives, with $m(s)=\int_{|v|<s}\chi(v)\,dv$,
\begin{equation}\label{eq:rate-newton}
 (\chi*U_D)(y)=\frac{2\kappa}{\omega_d}\int_D
 \nabla\rho(v)\cdot\frac{v-y}{|v-y|^d}m(|v-y|)\,dv.
\end{equation}
Indeed the spherical average of $(v-y-w)/|v-y-w|^d$ on $|w|=s$
is $(v-y)|v-y|^{-d}\mathbf1_{\{|v-y|>s\}}$; integrate in $s$.
Since the contact integrand is nonnegative,
$(\chi*U_E)(y_0)\le\|\chi\|_1H$.

In the plane take $\chi(v)=(1+|v|)^{-2}\mathbf1_{\{|v|\le Cr\}}$,
with the cutoff large enough to cover $D_n-y_0$.
Cell comparison and \eqref{eq:rate-crude} give
\[
 W_n(z)\le C(\chi*\ell_n)(y_0)
 \le C(\chi*U_{bK})(y_0)+C(H+\delta)\|\chi\|_1.
\]
The cone $U_{bK}=2d\kappa(b-\rho)_+$ is Lipschitz and vanishes at
$y_0$, so its convolution at $y_0$ is $O(r)$. Also
$\|\chi\|_1=O(a)$ and $H\le\delta$ by \eqref{eq:rate-crude}.
Consequently $W_n(z)\le C(r+\sqrt r\,a^2)\le Cr$, proving the claim.

For $d\ge3$ put $\chi(v)=(1+|v|)^{2-2d}$. This kernel is integrable
and $\chi*\chi\le A\chi$: split the convolution into
$|w|\ge|v|/2$ and $|v-w|\ge|v|/2$. Cell comparison, the martingale
bound and Young's inequality yield, for every fixed $\gamma>0$,
\[
 \ell_n\le U_{D_n}+C_\gamma a+\gamma\chi*\ell_n.
\]
This inequality holds globally; outside $B(0,2n)$ the local time is zero
and $|U_{D_n}|\le C$. Convolve it with $\chi$ and take
$\gamma=(2A)^{-1}$ to obtain the resolvent estimate
\begin{equation}\label{eq:rate-resolvent}
 \chi*\ell_n\le2\chi*U_{D_n}+Ca.
\end{equation}
The cone is Lipschitz, bounded by $Cb$, and $b\le Cr$, hence
\begin{equation}\label{eq:rate-cone-convolution}
 (\chi*U_{bK})(y)\le\|\chi\|_1 U_{bK}(y)+Ca;
\end{equation}
the error is bounded by $C\int\chi(v)\min\{|v|,r\}\,dv=O(a)$.
At $y_0$, \eqref{eq:rate-newton} and \eqref{eq:rate-resolvent} now give
$W_n(z)\le C(\chi*\ell_n)(y_0)\le C(H+a)$.
Substitute in \eqref{eq:rate-contact-noise} and absorb
$C\sqrt{aH}\le H/2+Ca$.
\end{proof}

\paragraph{Rates for every convex body.}
Lemma~\ref{lem:core-geometry} and the contact bound give
$I\le Cr^d h=Cr^{d+1}q$. Split $E$ into the layer
$b\le\rho<b+s$, whose volume is at most $Cr^{d-1}s$, and the remainder,
whose volume is at most $I/s$. With $s=r\sqrt q$ this proves
\begin{equation}\label{eq:rate-general-mass}
 |E|\le Cr^d\sqrt q,\qquad
 \|U_E\|_\infty\le Crq^{1/(d+1)}.
\end{equation}
The quadratic martingale is
\[
 Q_n=|X_n|^2-n+2\kappa\sum_{x\in A_n}\lambda(x)\rho(x).
\]
Stopping on leaving $B(0,Cr)$ and applying Azuma gives
$|Q_n|\le Cr\sqrt{na}$. The sparse-weight condition and cell quadrature
each cost $O(r^d)$ in replacing the weighted sum by $\int_{D_n}\rho$.
Since $\int_{bK}\rho=d|K|b^{d+1}/(d+1)$ and
$n=2d\kappa|K|r^{d+1}/(d+1)$, we obtain
\begin{equation}\label{eq:rate-clock-comparison}
 \left|1-(b/r)^{d+1}\right|
 \le C\left(r^{(1-d)/2}\sqrt a+r^{-1}+|E|r^{-d}\right).
\end{equation}
The first two terms are at most $C\sqrt q$; hence
\begin{equation}\label{eq:rate-general-conclusion}
 |b-r|\le Cr\sqrt q,\qquad
 \max_x|\ell_n(x)-2d\kappa(r-\rho(x))_+|
 \le Crq^{1/(d+1)}.
\end{equation}
Here we used the cone identity, \eqref{eq:rate-crude} and
\eqref{eq:rate-general-mass}; their errors are bounded by the right side.
Apply the projection crossing estimate \eqref{eq:core-crossing} with
$k=\lfloor r\rfloor$. Outside $kK$ the local time is at most
$C(1+rq^{1/(d+1)})$. Since $\rho$ is Lipschitz, it follows that
\begin{equation}\label{eq:rate-general-outer}
 \max_{j\le n}\rho(X_j)\le r+Crq^{1/(d+1)}a.
\end{equation}
These prove Theorem~\ref{thm:main}\textup{(ii)}. Taking $p>1$ and
applying Borel--Cantelli gives the filled inclusions and uniform cone
law in part \textup{(i)}; at each fixed site the cone grows linearly in
$r_n$, proving recurrence.

\section{Euclidean geometry and exterior probes}
Now let $\rho(v)=|v|$. We retain $r,a,b,E,h,q$ above, so $b=R_{\rm in}(n)$.
For $|v|\ge b=|y_0|$ the exact identity
\[
 2(|v|^2-v\cdot y_0)=|v-y_0|^2+|v|^2-b^2
\]
gives
\[
 u_v\cdot\frac{v-y_0}{|v-y_0|^d}
 \ge\frac{|v-y_0|^{2-d}}{2|v|}\ge c r^{1-d}
 \qquad(v\in E).
\]
Thus the contact estimate improves the volume and inner radius bounds to
\begin{equation}\label{eq:rate-euclidean-inner}
 |E|\le Cr^{d-1}h=Cr^dq,\quad |b-r|\le Crq,
 \quad |D_n\mathbin\triangle B(0,r)|\le Cr^dq.
\end{equation}
The radius assertion follows from \eqref{eq:rate-clock-comparison} since
its first two terms are at most $Cq$.

Write $U_E^+$ for the radial potential with its negative integrand
replaced by zero, and put
\[
 L=1+\max_{|x|\ge b}\ell_n(x),\qquad
 w=R_{\rm out}(n)-b+2\sqrt d,\qquad R=b+w.
\]
Then $w\ge1$ and $E\subset\{b\le|v|\le R\}$.
Crossing with $k=\lceil b\rceil$ gives $w\le CLa$.
For $d\ge3$, \eqref{eq:rate-newton} gives
\[
 \chi*U_E\le\|\chi\|_1U_E^+,\qquad
 \chi*U_E\le\|\chi\|_1\|U_E\|_\infty.
\]
Combining these with the resolvent and cone-convolution estimates yields
\begin{equation}\label{eq:rate-envelope}
 \ell_n(y)\le C\bigl((b-|y|)_++\|U_E\|_\infty+a\bigr),
 \quad
 \ell_n(y)\le C(U_E^+(y)+a)\quad(|y|\ge b).
\end{equation}
In the plane use \eqref{eq:rate-crude} instead. In either case the
volume bound $\|U_E\|_\infty\le C|E|^{1/d}$ gives
\begin{equation}\label{eq:rate-preliminary-width}
 L\le Crq^{1/d},\qquad w\le Crq^{1/d}a=o(r a^{-d}).
\end{equation}
In particular $b\sim r$ and $w=o(r)$. The layer estimate also gives
$\|U_E\|_\infty\le Cw$.

\paragraph{One exterior estimate in all dimensions.}
Set
\begin{equation}\label{eq:rate-theta}
 \Theta=\begin{cases}
 h,&d=2,\\
 a+r^{-1}\displaystyle\sup_y\int_E|v-y|^{2-d}\,dv,&d\ge3.
 \end{cases}
\end{equation}
In either case $\Theta\ge a$ for large $n$ and
$\sup_y\int_E|v-y|^{2-d}\,dv\le Cr\Theta$; in the plane this is
just $|E|\le Crh$. For $|y|\ge R$, the Euclidean identity therefore gives
$U_E^+(y)\le Cr^{-1}\int_E|v-y|^{2-d}\,dv\le C\Theta$.

The planar substitute for the resolvent is particularly short.
The layer estimate and \eqref{eq:rate-crude} give
$\ell_n(x)\le4\kappa(b-|x|)_++C(w+\sqrt r\,a)$.
For every lattice $y$ with $|y|\ge b$, inserting this in $W_n(y)$ gives
\[
 W_n(y)\le Cr+C(w+\sqrt r\,a)a\le Cr.
\]
Indeed $(b-|x|)_+\le|x-y|$ bounds the cone contribution by $Cr$,
and \eqref{eq:rate-preliminary-width} makes the other term $o(r)$.
Thus the martingale bound, and decay beyond the pole set, imply
\begin{equation}\label{eq:rate-planar-exterior}
 |\ell_n(y)-U_E(y)|\le Ch\qquad(d=2,\ |y|\ge b).
\end{equation}
Together with \eqref{eq:rate-envelope}, this gives in every dimension
$\ell_n(y)\le C(U_E^+(y)+\Theta)$ outside the inner ball.

We now use the same exterior probe in every dimension. Given a unit
vector $u$ and $\sqrt d\le t\le r/4$, put $y_t=(R+4t)u$ and let $z$
be its cell center. Since $R=R_{\rm out}+2\sqrt d$, this cell is empty
and $|z|\ge R$. If $d\ge3$, the resolvent, cone-convolution estimate
and $\chi*U_E\le\|\chi\|_1U_E^+$ give $W_n(z)\le C\Theta$;
if $d=2$, use \eqref{eq:rate-planar-exterior}.
The martingale estimate and cell modulus give in both cases
\begin{equation}\label{eq:rate-probe}
 |U_E(y_t)|\le C\Theta.
\end{equation}
For $v\in E\cap B(Ru,t)$ we have $u_v\cdot u\ge1-t/R$ and
\[
 u_v\cdot(v-y_t)\le R-(R+4t)(1-t/R)\le-t,
 \qquad |v-y_t|\le5t,
\]
because $t\le R/2$ for large $n$. Such points contribute at most
$-ct^{1-d}$ each, whereas all positive contributions total at most
$C\Theta$. Hence \eqref{eq:rate-probe} implies
\begin{equation}\label{eq:rate-cap}
 |E\cap B(Ru,t)|\le C\Theta t^{d-1}\qquad(|u|=1,\ t\ge1).
\end{equation}
For $t>r/4$ use $|E|\le Cr^{d-1}h\le Cr^{d-1}\Theta$;
for $1\le t<\sqrt d$ use the ball-volume bound. Thus the estimate
holds down to a fixed spatial scale, independently of the shell width.

\begin{lemma}[Cap volume controls the boundary profile]\label{lem:rate-near-far}
For every $y$ with $b\le|y|\le R$, put $s=R-|y|+1$. Then
\[
 U_E^+(y)\le C(\Theta s^{d-1})^{1/d}+C\Theta.
\]
The same bound holds for $\ell_n(y)$ at lattice sites.
\end{lemma}
\begin{proof}
For $t\ge s$, $B(y,t)\subset B(Ru_y,2t)$, so
\eqref{eq:rate-cap} gives $|E\cap B(y,t)|\le C\Theta t^{d-1}$.
The part $|v-y|\le s$ has volume at most $C\Theta s^{d-1}$ and
therefore potential at most $C(\Theta s^{d-1})^{1/d}$.
For the rest, $|v|^2-|y|^2\le2Rs$ and $|v|\ge b\asymp r$ give
\[
 \left[u_v\cdot\frac{v-y}{|v-y|^d}\right]_+
 \le\frac{Cs}{|v-y|^d}+\frac{C}{r|v-y|^{d-2}}.
\]
The cap bound makes the first integral at most
$Cs\int_s^\infty\Theta t^{d-1}t^{-d-1}\,dt\le C\Theta$;
the definition of $\Theta$ bounds the second by $C\Theta$.
The exterior local-time bound proves the last assertion.
\end{proof}

Since $s\le w+1\le2w$, the lemma and crossing give
\begin{equation}\label{eq:rate-L}
 L\le C\bigl((\Theta w^{d-1})^{1/d}+\Theta\bigr),
 \qquad w\le CLa\le\tfrac12w+C\Theta a^d.
\end{equation}
Thus $w\le C\Theta a^d$ in every dimension. In the plane $\Theta=h$,
so the lemma immediately yields
\begin{equation}\label{eq:rate-planar-final}
 w\le C\sqrt r\,a^{5/2},\qquad
 \sup_{b\le|y|\le R}U_E^+(y)\le C\sqrt r\,a^{3/2}.
\end{equation}

For $d\ge3$ it remains only to absorb the curvature term in $\Theta$.
If $t\ge w$ and $B(y,t)$ meets $E$ at $v$, its intersection with $E$
lies in $B(v,2t)$ and hence in a cap of radius at most $3t$.
Consequently
\[
 |E\cap B(y,t)|\le C\min\{\Theta t^{d-1},r^{d-1}a\}\quad(t\ge w),
 \qquad |E\cap B(y,t)|\le\omega_dt^d\quad(t<w).
\]
Layer-cake integration, splitting the minimum at
$t=r(a/\Theta)^{1/(d-1)}$, gives
\[
 \sup_y\int_E|v-y|^{2-d}\,dv
 \le Cw^2+Cr\Theta^{(d-2)/(d-1)}a^{1/(d-1)}.
\]
Substitution in \eqref{eq:rate-theta} and Young's inequality give
\begin{equation}\label{eq:rate-theta-improved}
 \Theta\le C(a+w^2/r).
\end{equation}
Together with $w\le C\Theta a^d$, this implies
$w\le Ca^{d+1}+Ca^dw^2/r$. The last term is at most $w/2$ for
large $n$ by \eqref{eq:rate-preliminary-width}. We conclude that
\begin{equation}\label{eq:rate-higher-final}
 w\le Ca^{d+1},\qquad \Theta\le Ca.
\end{equation}
Substitution in \eqref{eq:rate-cap} proves the all-scale cap estimates.
Lemma~\ref{lem:rate-near-far} now also gives
\begin{equation}\label{eq:rate-boundary-profile}
 \ell_n(y)\le
 \begin{cases}
 C\bigl(h+\sqrt{hs}\bigr),&d=2,\\
 C\bigl(a+(as^{d-1})^{1/d}\bigr),&d\ge3,
 \end{cases}
 \qquad b\le|y|\le R,\quad s=R-|y|+1.
\end{equation}
In particular, for $d\ge3$ the outermost radial layer of width
$C_0\log n$ has local times $O(\log n)$. If this layer reaches inside
$B(0,b)$, then $w=O(a)$. The continuum boundary bound from the lemma
and the maximum principle give $U_E\le Ca$ inside that ball.
There $\chi*U_E\le\|\chi\|_1U_E$, since the potential integrand is
nonnegative; the resolvent and the cone depth $O(a)$ give the same
local-time bound.

\paragraph{Uniform and interior local times.}
For $d\ge3$, the layer bound and \eqref{eq:rate-higher-final} give
$\|U_E\|_\infty\le Ca^{d+1}$. Together with the cone identity,
\eqref{eq:rate-crude} and \eqref{eq:rate-euclidean-inner}, this gives
\[
 \max_x|\ell_n(x)-2d\kappa(r-|x|)_+|\le C\sqrt{ra}.
\]
In the plane $U_E$ is nonnegative and harmonic on $B(0,b)$, and continuous
on its closure. The maximum principle and \eqref{eq:rate-planar-final}
therefore bound it throughout that ball by $C\sqrt r\,a^{3/2}$.
On the outer annulus use \eqref{eq:rate-planar-exterior}; beyond the range
both the local time and, except for a radial error $C|b-r|$, the target
cone vanish. Consequently the uniform planar error is
$C\sqrt r\,a^{3/2}$.

Finally, for any fixed $\theta<1$ and $|y|\le\theta r$, the excess set
is at distance at least $c_\theta r$ for large $n$. Thus
\begin{equation}\label{eq:rate-bulk-potential}
 |U_{D_n}(y)-2d\kappa(r-|y|)|
 \le C_\theta(r^{1-d}|E|+|b-r|)\le C_\theta rq.
\end{equation}
The corresponding local-time error is at most $C_\theta\sqrt r\,a$
in the plane and $C_\theta\sqrt{ra}$ in higher dimensions.
All assertions hold simultaneously with probability $1-Cn^{-p}$ and,
by Borel--Cantelli, eventually almost surely. The all-scale cap bounds
\eqref{eq:rate-cap} and boundary profile \eqref{eq:rate-boundary-profile}
are additional quantitative conclusions of the same proof.
